\iffalse
Counterfactual explanations in classification demonstrate why a decision was made for option A instead of option B. In symbolic models, this modification can be significant and indicate causality between feature values and class \cite{Lewis1973}. However, in high-dimensional spaces where causal models are not available, this modification is reduced to an adversarial attack~\cite{Madry2017}. Adversarial attacks require carefully chosen small modifications, such as imperceptible noise, to change the class of an example and deceive the network. As a result, this noise does not typically provide genuine counterfactual explanation since it is counterfactual with respect to the model but not to the class \cite{wachter2017counterfactual}. Thus, Saliency Maps \cite{SimonyanVZ13} -- gradient of output with respect to the input --, which are the basis for most adversarial attacks, are generally unsuitable for explaining model decisions. Therefore, several methods requiring more complex computations, such as SmoothGrad \cite{Smilkov2017}, Integrated Gradient \cite{Sundararajan2017}, or Grad-CAM \cite{Selvaraju_2019}, have been proposed to provide better explanations. Recently, the XAI community has investigated the link between explainability and robustness and proposed methods and metrics accordingly~\cite{hsieh2020evaluations, boopathy2020proper, lin2019explanations, ross2021learning}.
In~\cite{serrurier2021achieving}, authors propose to address the weakness with respect to adversarial attacks by training \onlyonelip~constrained neural networks with a loss that is the dual of an optimal transport optimization problem, called hKR and noted \losshkr. The resulting models have been shown to be robust with a certifiable margin. We refer to these networks as Optimal Transport Neural Networks (\otnn) hereafter.
\fi


The Lipschitz constant of a function expresses the extent to which the output may vary for a small shift in the input. As a composition of numerous functions, the Lipschitz constant of a neural network can be arbitrarily high, particularly when trained for a classification task \cite{bethunelip}. Adversarial attacks~\cite{Madry2017} exploit this weakness by selecting minor modifications, such as imperceptible noise, for a given example to change the predicted class and deceive the network. Consequently, Saliency Maps \cite{SimonyanVZ13} -- gradient of output with respect to the input --, serve as the basis for most adversarial attacks and often highlight noisy patterns that fool the model instead of meaningful modifications, rendering them generally unsuitable for explaining model decisions. Therefore, several methods requiring more complex computations, such as SmoothGrad \cite{Smilkov2017}, Integrated Gradient \cite{Sundararajan2017}, or Grad-CAM \cite{Selvaraju_2019}, have been proposed to provide smoother explanations. Recently, the XAI community has investigated the link between explainability and robustness and proposed methods and metrics accordingly~\cite{hsieh2020evaluations, boopathy2020proper, lin2019explanations, ross2021learning}.
%\commentgreen{Such  automatic metrics can still suffer from biases, and no clear evidence that they are correlated with human notion of explanations. \cite{fel2022aligning} suggests comparing the alignment between attribution methods and human feature importance using the ClickMe dataset~\cite{linsley2017visual}.}
However, the reliability of those automatic metrics can be compromised by artifacts introduced by the baselines ~\cite{hsieh2020evaluations,sturmfels2020visualizing,haug2021baselines,kindermans2019reliability,hase2021out}, and there is no conclusive evidence demonstrating their correlation with the human understanding of explanations. To address this, a study by \cite{fel2022aligning} suggests completing those metrics with the alignment between attribution methods and human feature importance using the ClickMe dataset~\cite{linsley2017visual}.
%\mathieu{To obtain a metric that better mirrors human interpretation of decisions, \cite{fel2022aligning} suggests comparing the alignment between attribution methods and human feature importance using the ClickMe dataset~\cite{linsley2017visual}.}

In~\cite{serrurier2021achieving}, authors propose to address the weakness with respect to adversarial attacks by training \onlyonelip~constrained neural networks with a loss that is the dual of an optimal transport optimization problem, called hKR  (Eq.~\ref{eq:reg_OT}). The resulting models have been shown to be robust with a certifiable margin. We refer to these networks as Optimal Transport Neural Networks (\otnn) hereafter.

In this paper, we demonstrate that OTNNs exhibit valuable explainability properties. Our experiments reveal that OTNN Saliency Maps significantly outperform various attribution methods for unconstrained networks across all tested state-of-the-art Explainable XAI metrics. This improvement is consistent across toy datasets, large image datasets, binary, and multiclass problems. Qualitatively, OTNN Saliency Maps concentrate on crucial image regions and generate less noise than maps of unconstrained networks, as illustrated in Figure \ref{fig:big_picture}. Figure \ref{fig:big_picture}.c presents the Saliency Maps of an OTNN based on ResNet50 alongside its unconstrained vanilla counterpart, both trained on ImageNet~\cite{imagenet_cvpr09}. In the unconstrained case, the Saliency Maps appear sparse and uninformative, with the most critical pixels often located outside the subject. Conversely, the OTNN Saliency Map is less noisy and highlights significant image features. This distinction is emphasized in Figure \ref{fig:big_picture}.d), comparing the feature visualization of the two models. Feature visualization extracts the inverse prototypical image for a given class using gradient ascent \cite{olah2017feature,nguyen2016multifaceted}. The results for vanilla ResNet are noticeably noisy, making class identification difficult. In contrast, feature visualization with OTNN yields clearer results, displaying easily identifiable features and classes (e.g., goldfish, butterfly, and medusa). Furthermore, modifying the image following the gradient direction provides an interpretable method for altering the image to change its class. Pictures \ref{fig:big_picture}.a) and \ref{fig:big_picture}.b) display the original image, the gradient direction, and the transformation following the gradient direction (refer to Section 4 for details). We observe explicit structural changes in the images, transforming them into an example of another class in both multiclass (MNIST) and high-resolution multi-labels cases (e.g., smile/not smile and blond hair/not blond hair classification). Lastly, a large-scale human experiment demonstrates that these maps are remarkably aligned with human attribution on ImageNet (Fig. \ref{fig:human_alignement}).

%Fig. \ref{fig:mnist} illustrates, on an \otnn~learned on MNIST dataset, the transformation of an image of the class zero with respect to the gradient of another class's output (as done in a vanilla targeted attack). We observe that it explicitly changes the zero into the target number, and gradients provide understandable explanations of why it was not this number.

\begin{figure*}[ht]
  \centering
  %\fbox{\includegraphics[width=0.95\linewidth]{img/mnist.png}}
  %\includegraphics[width=0.99\linewidth]{img/big_picture.jpg}
  \includegraphics[width=0.99\linewidth]{img/big_picturev_2.pdf}
\caption{
\textbf{Illustration of the beneficial properties of OTNN gradients.} Examples \textbf{a)} and \textbf{b)} show that the gradients naturally provide a direction that enables the generation of adversarial images - a theoretical justification based on optimal transport is provided in Section 3. By applying the gradient $\x' = \x  -t \nabla_{\x}  \f(\x)$ to the original image $\x$ (on the left), any digit from MNIST can be transformed into its counterfactual $\x'$ (e.g., turning a 0 into a 5). In \textbf{b)}, we illustrate that this approach can be applied to larger datasets, such as Celeb-A, by creating two counterfactual examples for the closed-mouth and blonde  classes. In \textbf{c)}, we compare the 
Saliency Map
%natural gradient 
of a classical model with those of OTNN gradients, which are more focused on relevant elements. Finally, in \textbf{d)}, we show that following the gradients of OTNN could generate convincing feature visualizations that ease the understanding of the model's features. 
}
\label{fig:big_picture}
\end{figure*}

%We provide a theoretical justification of the the well behavior of OTNN salency map. Starting from the fact that \otnn s~encode the dual formulation of the optimal transport problem, and we prove that the gradient of a given point $x$ is both (i) in the direction of the closest adversarial example on the decision boundary and (ii) in the direction of the image of $x$ according to the underlying transport plan. This implies that adversarial attacks for an \otnn~are equivalent to travels along the optimal transport path. As a result, the modification obtained is not only an adversarial attack but also a counterfactual explanation, i.e., an explanation of why the classification was not another class\cite{Lewis1973}. An optimal transport plan between two classes can be then interpreted as a global approach for constructing counterfactualsas suggested in \cite{black19OptimTransp,deLara2021}. These counterfactuals no longer necessarily correspond to the smallest transformation for a given input sample, but rather the smallest transformation on average when pairing points from two classes. The consequence of this property is that the Saliency Map of \otnn s for an image indicates the importance of each pixel in the modifications required to change the class. It should be noted that several methods based on GAN \cite{Jacob2021} or causality penalty \cite{karimi2021algorithmic} achieve highly realistic counterfactual images. In this paper, we do not aim to compete with the quality of these results but to demonstrate that \otnn s Saliency Maps possess theoretically and empirically properties of counterfactual explanations.

We provide a theoretical justification for the well-behaved OTNN Saliency Maps. Building upon the fact that OTNNs encode the dual formulation of the optimal transport problem, we prove that the gradient of the optimal solution at a given point $\x$ is both (i) in the direction of the nearest adversarial example on the decision boundary, and (ii) in the direction of the image of $x$ according to the underlying transport plan. This implies that adversarial attacks for an OTNN are equivalent to traversing the optimal transport path which can be achieved  by following the gradient. Consequently, the resulting modification serves both as an adversarial attack and a counterfactual explanation, 
%which explains why the classification was not another class
explaining why the decision was A and not B
\cite{Lewis1973}. An optimal transport plan between two classes can be interpreted as a global approach for constructing counterfactuals, as suggested in \cite{black19OptimTransp,deLara2021}. These counterfactuals may not correspond to the smallest transformation for a given input sample but rather the smallest transformation on average when pairing points from two classes. A consequence of this property is that the Saliency Map of an OTNN for an image indicates the importance of each pixel in the modifications needed to change the class. It is worth noting that several methods based on GAN \cite{Jacob2021} or causality penalty \cite{karimi2021algorithmic} produce highly realistic counterfactual images. %In this paper, we do not seek to compete with the quality of these results but to demonstrate that OTNN Saliency Maps possess theoretically and empirically grounded properties of counterfactual explanations. 
However, the objective of our paper is not to compete with the quality of these results, but rather to demonstrate that OTNN Saliency Maps possess both theoretical and empirical foundations as counterfactual explanations.

%\commentgreen{EN LISANT TOUT CELA JE ME DEMANDE SI ON EN METS PAS TROP DANS l'INTRO: on affiche tous les arguments du papier même la causalité c'est plus une intro mais un résumé long? Vos avis?}
%Second, we propose a method to suggest an enhancements for the multiclass version of the hKR loss proposed in \cite{serrurier2021achieving}, resulting in improved performance even on large scale problems such as ImageNet. 

We summarize our contributions as follows: first, after introducing the background on OTNN and XAI, we establish several properties of the gradient of an OTNN with respect to adversarial attacks, decision boundaries, and optimal transport. Second, we establish that the optimal transport properties of OTNN's gradient lead to a reinterpretation of adversarial attacks as counterfactual explanations, consequently endowing the Saliency Map with the favorable XAI properties inherent in these explanations. Third, our experiments support the theoretical results, showing that 
%Saliency Maps significantly improve metric scores for most XAI methods when compared to their use on unconstrained neural networks. 
%Additionally, we find that the simple Saliency Map of an OTNN achieves top-ranked scores on state-of-the-art XAI metrics compared to more sophisticated methods and is equivalent to those provided by Smoothgrad. 
metric scores are higher for most of the XAI methods on \otnn~compared to unconstrained neural networks. Additionally, we find that the Saliency Map for \otnn~achieves top-ranked  scores on XAI metrics compared to more sophisticated XAI methods, and is equivalent to Smoothgrad. 
Lastly, drawing from \cite{fel2022aligning}, we emphasize that OTNNs are naturally and remarkably aligned with human explanations, and we present several examples of gradient-based counterfactuals obtained with OTNNs. 


